% SEGMENT OLP-0717-S01
% Part: sets-functions-relations
% Chapter: functions
% Section: isomorphisms

\documentclass[../../../include/open-logic-section]{subfiles}


\begin{document}

\olfileid{sfr}{fun}{iso}
\olsection{ஓரினச்சார்பு}

\begin{explain}
\emph{ஓரினச்சார்பு} என்பது, அது தொடர்புபடுத்தும் கணங்களின்
அமைப்பைப் பேணும் ஓர் இருபுறச் சார்பு; இங்கு அமைப்பு என்பது
அக்கணங்களின் !!{element}s இடையே நிலவும் உறவுகளைச் சார்ந்தது.
பின்வரும் இரண்டு கணங்களைக் கருதுக: $X=\{1,2,3\}$ மற்றும்
$Y=\{4,5,6\}$. இரு கணங்களுக்கும் தொடரி, சிறியது,
பெரியது ஆகிய உறவுகளால் அமைப்பு வழங்கப்படுகிறது.
இரு கணங்களுக்கிடையிலான ஓரினச்சார்பு அவ்வமைப்புகளைப் பேணும்
!!a{bijection} ஆகும். எனவே !!a{bijective} சார்பு $f \colon X
\to Y$ ஆனது $i<j$ அப்போதும் அப்போது மட்டுமே $f(i)<f(j)$,
$i>j$ அப்போதும் அப்போது மட்டுமே $f(i)>f(j)$, மேலும் $j$ ஆனது
$i$ இன் தொடரி அப்போதும் அப்போது மட்டுமே $f(j)$ ஆனது
$f(i)$ இன் தொடரி எனில், அது ஓரினச்சார்பு ஆகும்.
\end{explain}

% SEGMENT OLP-0717-S02
\begin{defn}[ஓரினச்சார்பு]
$U$ என்பது $\langle X, R\rangle$ என்ற ஜோடியாகவும், $V$ என்பது
$\langle Y, S\rangle$ என்ற ஜோடியாகவும் இருக்கட்டும்; இங்கு
$X$, $Y$ கணங்கள், $R$, $S$ முறையே $X$, $Y$ மீதான உறவுகள்.
$X$ இலிருந்து $Y$ க்குச் செல்லும் !!a{bijection}~$f$ உறவு
அமைப்பைப் பேணினால், அப்போதும் அப்போது மட்டுமே, அது $U$ இலிருந்து
$V$ க்குச் செல்லும் \emph{ஓரினச்சார்பு} ஆகும். அதாவது,
$X$ இல் உள்ள எந்த $x_{1}$, $x_{2}$ க்கும்
$\tuple{x_1,x_2} \in R$ அப்போதும் அப்போது மட்டுமே
$\tuple{f(x_1),f(x_2)} \in S$.
\end{defn}

% SEGMENT OLP-0717-S03
\begin{ex}
பின்வரும் இரண்டு கணங்கள் $X=\{1,2,3\}$,
$Y=\{4,5,6\}$ மற்றும் சிறியது, பெரியது என்ற உறவுகளைக்
கருதுக. $f(x) = 7-x$ என வரையறுக்கப்படும்
$f\colon X \to Y$ என்ற சார்பு, $\tuple{X,<}$ க்கும்
$\tuple{Y,>}$ க்கும் இடையிலான ஓரினச்சார்பு ஆகும்.
\end{ex}

\end{document}
